\section{Background \& Threat Model}
\label{sec:background}

\subsection{Programmatic Advertising \& Header Bidding}
Modern web publishers monetize digital inventory via programmatic Real-Time Bidding (RTB)~\cite{castelluccia2014selling, bashir2016tracing}. 
While legacy display monetization relied on monolithic waterfall ad servers (e.g., Google Ad Manager), modern publishers deploy \textbf{Header Bidding} wrappers, principally the open-source \texttt{Prebid.js} framework, embedded directly in the client-side DOM to conduct parallel auctions prior to ad-server calls~\cite{pachilakis2019no, cook2020inferring}.

When a client navigates to a publisher, the programmatic auction resolves through four sequential stages~\cite{pachilakis2019no, barysei2023investigating}:
\begin{enumerate}
    \item \texttt{Prebid.js} reads configuration tags, identifies active ad slots, and dispatches parallel, asynchronous HTTP bid requests to configured Supply-Side Platform (SSP) adapters (e.g., Rubicon, OpenX, AppNexus)~\cite{cook2020inferring};
    \item Each SSP adapter forwards the bid request via OpenRTB protocols to participating Demand-Side Platforms (DSPs), encapsulating client attributes: geographic IP localization, viewport dimensions, historical cookie-synchronization tokens~\cite{papadopoulos2019cookie}, smuggled navigational identifiers~\cite{randall2022measuring}, and stateless device fingerprint hashes extracted by embedded scripts~\cite{papadogiannakis2021user};
    \item DSP bidding algorithms parse the payload and query internal identity graphs: matching an affluent persona triggers aggressive premium bidding (\$2.50--\$15.00 CPM), whereas unresolved identities retreat to generic contextual reserve floors (\$0.05--\$0.25 CPM)~\cite{barysei2023investigating};
    \item Winning bids return to the browser, where \texttt{Prebid.js} resolves the highest clearing \cpm, surfaces it to the publisher's ad server, and renders the winning creative inside an isolated \texttt{iframe}.
\end{enumerate}

\subsection{The Economic Targeting Premium}
\label{subsec:targeting_premium}

The foundation of programmatic digital advertising is that identified users command significantly higher clearing prices than un-profiled users~\cite{castelluccia2014selling, bashir2016tracing}. Longitudinal empirical analysis by Beales and Eisenach established that behavioral-targeted advertising sells for a $2.68\times$ price premium over non-targeted contextual advertising ($+168\%$)~\cite{beales2014empirical}. 

Subsequent field trials by Google confirmed that stripping user identifiers decreases publisher ad revenue by an average of $52\%$~\cite{google2019cookie}, an empirical finding corroborated by Johnson et al., who documented an identical $52\%$ ad revenue loss when consumers opt out of behavioral profiling~\cite{johnson2020consumer}. 

Consequently, we formalize auction clearing \cpm as an \textit{economic side-channel oracle}:
\begin{equation}
\cpm(u) = \alpha \cdot \mathbb{I}_{\text{identified}}(u) + \beta_{\text{context}} + \epsilon_{\text{market}}
\label{eq:economic_oracle}
\end{equation}
where $\alpha$ represents the behavioral targeting premium ($\alpha > 0$)~\cite{beales2014empirical, johnson2020consumer}, $\beta_{\text{context}}$ represents baseline contextual inventory value, and $\epsilon_{\text{market}}$ denotes stochastic auction noise driven by DSP pacing algorithms and exchange latency~\cite{barysei2023investigating}.

\subsection{Threat Model}
\label{subsec:threat_model}

\paragraph{Adversarial Capabilities}
We model commercial tracking consortiums, SSPs, and DSPs under the classical \textbf{honest-but-curious} (semi-honest) adversarial model operating across the programmatic advertising mesh~\cite{papadopoulos2019cookie, oh2022cart}. 
The adversary operates within standard unprivileged client execution contexts, possessing legitimate access to standard DOM and Web APIs~\cite{bujlow2017survey, laperdrix2020browser}. 
Embedded third-party scripts participate in header bidding auctions and attempt to persist user identities across disjoint first-party administrative domains ($s_i \neq s_j$) without user consent~\cite{papadogiannakis2021user}. 

The adversary leverages standard stateless fingerprinting vectors documented in the literature: Canvas pixel extraction~\cite{mowery2012pixel, acar2014web}, WebGL parameter enumeration~\cite{cao2017cross}, AudioContext oscillator decay~\cite{englehardt2016openwpm, laperdrix2020browser}, clock skew timing attacks~\cite{sanchez2018clock}, browser extension profiling~\cite{starov2019unnecessarily}, and WebRTC candidate harvesting.

\paragraph{Adversarial Boundaries}
Following the canonical web tracking threat model established in the literature~\cite{tor_design, englehardt2016openwpm, laperdrix2020browser}, the adversary is strictly bounded by the browser client-side execution sandbox and standard web protocol guarantees. 
Specifically, the adversary \textit{cannot} compromise the underlying operating system kernel, inspect browser process memory, forge TLS cryptographic primitives, or mount active machine-in-the-middle network tampering~\cite{bujlow2017survey}. 
Coarse IP-based location clustering is constrained by dynamic residential subnet reassignment and carrier-grade NAT variance~\cite{eckersley2010unique, englehardt2016openwpm}. 
Finally, server-side identity synthesis entirely detached from client-observable hardware signals (e.g., first-party authenticated CRM matching) is orthogonal to client-side browser defense efficacy~\cite{bashir2016tracing, cook2020inferring}.

\paragraph{Client Defenses}
The client executes a modern web browser equipped with anti-fingerprinting defenses:
\begin{enumerate}
    \item \textbf{API Normalization:} Tor Browser~\cite{tor_design, torbrowser_normalization_2019} and Firefox RFP~\cite{mozilla_rfp_docs}, which homogenize hardware attributes to standardized anonymity sets;
    \item \textbf{API Randomization:} Brave Farbling~\cite{brave_fingerprinting_protections_2025}, which injects cryptographic noise into Web API returns;
    \item \textbf{Defensive Network Filtering:} Request-level blocking via tracking protection lists and graph-based filters~\cite{iqbal2020adgraph}.
\end{enumerate}

The user's objective is to enforce an \textit{identity break} by programmatically flushing all stateful storage and verifying whether stateless defenses prevent commercial DSPs from re-linking the profile. 
However, capturing this economic signal requires an evasion-resilient measurement platform, as detailed next.